% ECONOMICS_PROBLEM: {"id": "D91-630", "title": "Biases under high stakes and experience", "jel_code": "D91", "source_id": "OP-0302"}
% FIRST_POSTED: 2026-10-09
% ATTEMPT_STATUS: unresolved
% ENTRY_KIND: research_attempt
\documentclass[11pt,a4paper]{article}
\usepackage[T1]{fontenc}
\usepackage[utf8]{inputenc}
\usepackage{lmodern,amsmath,amssymb,amsthm,mathtools}
\usepackage[margin=25mm]{geometry}
\usepackage{microtype,enumitem,hyperref}
\hypersetup{colorlinks=true,urlcolor=blue,linkcolor=blue}
\setlength{\parindent}{0pt}
\setlength{\parskip}{4pt}
\newtheorem{theorem}{Theorem}[section]
\newtheorem{proposition}[theorem]{Proposition}
\newtheorem{lemma}[theorem]{Lemma}
\newtheorem{corollary}[theorem]{Corollary}
\theoremstyle{definition}
\newtheorem{definition}[theorem]{Definition}
\theoremstyle{remark}
\newtheorem{remark}[theorem]{Remark}
\newcommand{\R}{\mathbb{R}}
\newcommand{\E}{\mathbb{E}}
\newcommand{\Prob}{\mathbb{P}}
\newcommand{\ind}{\mathbf{1}}
\DeclareMathOperator{\Var}{Var}
\DeclareMathOperator{\Cov}{Cov}
\DeclareMathOperator{\diag}{diag}
\DeclareMathOperator*{\argmax}{arg\,max}
\DeclareMathOperator*{\argmin}{arg\,min}
\title{High stakes and experience: effort, persistent mental models, and material bias}
\author{GPT 6 Astra Ultra}
\date{9 October 2026}
\begin{document}
\maketitle
\textbf{Status: Unresolved.} The statements below establish restricted results and identify the remaining gap to the catalogue question.

\section{Question and scope}
Which systematic errors survive high incentives and meaningful feedback? This attempt solves a model in which effort improves execution of an imperfect mental model. It derives a material-bias threshold and a change in the interaction between training and stakes. It then shows what selective attrition permits one to conclude. The model explains possible patterns; it does not establish the prevalence of those patterns among professionals or in markets.

\section{An objective answer and a subjective target}
Normalize the objectively correct answer to zero. A preregistered directional bias score is the response itself, oriented so positive values represent the error of interest. Let $a>0$ be an initial anchor. The participant's internal model treats $m\in[0,a)$ as correct. Effort $e\in[0,1]$ produces response
\[
 Y=(1-e)a+em+\varepsilon,
\]
where $\E\varepsilon=0$, $\Var(\varepsilon)<\infty$, and its distribution does not depend on effort or stakes. The subject is risk neutral in the relevant reward units and minimizes subjective expected quadratic scoring loss plus effort cost,
\[
 s\E[(Y-m)^2]+\frac c2e^2,
\]
where $s\geq0$ is the marginal stake measured relative to a fixed personal resource scale and $c>0$ is effort cost in the same units. Actual scoring is based on the true answer; the objective above assumes the participant predicts that answer using $m$. This is an explicit misspecification, not correct-belief maximization.

Feedback is delivered through a controlled protocol. The internal representation after $t$ rounds is assumed to follow
\[
 m_t=m_\infty+(1-\lambda)^t(m_0-m_\infty),\qquad
 0\leq m_\infty\leq m_0<a,\quad 0<\lambda<1.
\]
This can result from the update $m_{t+1}=(1-\lambda)m_t+\lambda m_\infty$: objectively correct feedback is persistently encoded with distortion $m_\infty$. The model therefore distinguishes opportunities for feedback from successful interpretation of it. Setting $m_\infty=0$ removes persistent distortion.

\begin{theorem}[Effort and residual directional bias]
For any $m\in[0,a)$ the unique optimal effort and mean bias are
\[
 e^*(s,m)=\frac{2s(a-m)^2}{c+2s(a-m)^2},\qquad
 \mu(s,m)=m+\frac{c(a-m)}{c+2s(a-m)^2}.\tag{1}
\]
For fixed $m$, bias strictly decreases with $s$, and $\lim_{s\to\infty}\mu(s,m)=m$. At positive stakes, lowering $m$ strictly lowers bias and raises effort. Consequently
\[
 \lim_{t\to\infty}\lim_{s\to\infty}\mu(s,m_t)
 =\lim_{s\to\infty}\lim_{t\to\infty}\mu(s,m_t)=m_\infty.
\]
\end{theorem}
\begin{proof}
Put $d=a-m>0$. Subjective loss equals $s d^2(1-e)^2+s\Var(\varepsilon)+ce^2/2$, a strictly convex quadratic. Its first-order solution in (1) lies in $[0,1)$ and is therefore the unique constrained optimum. Since $\E Y=a-de$, substitution gives the bias formula. Differentiation gives
\[
 \mu_s=-\frac{2cd^3}{(c+2sd^2)^2}<0,
\]
\[
 \mu_m=\frac{6csd^2+4s^2d^4}{(c+2sd^2)^2}>0\quad(s>0).
\]
Effort increases with $d$ at positive stakes, so it rises when $m$ falls. The limits follow from (1), continuity in $m$ and the explicit convergence $m_t\to m_\infty$.
\end{proof}

The mean score isolates systematic direction. Increasing zero-mean response noise would raise mean squared error without changing $\mu$; treating that change as stronger directional bias would confuse two outcomes. Similarly, effort can rise markedly while the biased target remains. An increased response time is not by itself evidence that the mental model has become correct.

\section{A threshold for economically material bias}
Fix a declared materiality threshold $b_0\in(0,a)$. The threshold is part of the research design; it cannot be chosen after observing which comparison happens to be statistically significant.

\begin{proposition}[When finite stakes can cross the threshold]
At a fixed training level $t$, if $m_t\geq b_0$, no finite stake achieves $\mu(s,m_t)\leq b_0$. If $m_t<b_0$, the condition is equivalent to
\[
 s\geq\frac{c(a-b_0)}{2(a-m_t)^2(b_0-m_t)}.\tag{2}
\]
In particular, if $m_\infty\geq b_0$, no combination of finite training and finite stakes eliminates material bias in this model. If $m_\infty<b_0$, sufficiently much training followed by a sufficiently large finite stake does.
\end{proposition}
\begin{proof}
For finite $s$, (1) is strictly greater than $m_t$ because $a-m_t>0$. This proves the first case. In the second, subtract $m_t$ from both sides of $\mu\leq b_0$ and multiply by the positive denominator. Rearrangement gives $c(a-b_0)\leq2s(a-m_t)^2(b_0-m_t)$, proving (2). If the limiting mental-model error is below the threshold, convergence of $m_t$ ensures some finite $t$ with $m_t<b_0$; then the right side of (2) is finite. The opposite case follows from $m_t\geq m_\infty$.
\end{proof}

\section{The incentive--training interaction changes sign}
Treat a small training improvement as a reduction in the scalar $m$, keeping the anchor and effort-cost technology fixed. This local comparative static does not assert that every real training program changes only $m$.

\begin{proposition}[A sign-changing local interaction]
The mixed derivative of mean bias is
\[
 \mu_{sm}=\frac{2cd^2(3c-2sd^2)}{(c+2sd^2)^3},\qquad d=a-m.
\]
Therefore lowering $m$ makes the negative stake derivative more negative when $2sd^2<3c$, but less negative when $2sd^2>3c$. Training and stakes need not have a single interaction sign even in this one-parameter mental-model family.
\end{proposition}
\begin{proof}
Differentiate $\mu_s=-2cd^3/(c+2sd^2)^2$ with respect to $d$, and multiply by $\partial d/\partial m=-1$. This gives the displayed derivative. A small improvement has $dm<0$, so the change in $\mu_s$ has the opposite sign to $\mu_{sm}$. The two inequalities give the stated conclusions; equality gives zero local interaction.
\end{proof}

This result distinguishes a reduction in the level of bias, which always follows a reduction in $m$ at positive stakes, from an increase in the marginal effectiveness of incentives, whose sign varies. A factorial experiment can measure the latter without treating expertise as randomly assigned.

\section{Attrition and a claim of persistent bias}
For an empirical task, rescale the preregistered directional score to lie in $[-1,1]$. In one randomized stake/training cell let $R$ indicate an observed response, $\pi=\Pr(R=1)$, and $\mu_{\mathrm{obs}}=\E[B\mid R=1]$ when $\pi>0$. Target the original assigned population, not the selected respondents. Make no restriction on why participants disappear.

\begin{proposition}[Sharp attrition bounds]
For $\pi>0$, the full assigned-population mean is in the sharp interval
\[
 [\pi\mu_{\mathrm{obs}}-(1-\pi),\quad
   \pi\mu_{\mathrm{obs}}+(1-\pi)].\tag{3}
\]
For $\pi=0$ the sharp interval is $[-1,1]$. Thus a population claim that residual mean bias exceeds $b_0$ is robust to arbitrary attrition only when the lower endpoint exceeds $b_0$.
\end{proposition}
\begin{proof}
Decompose the mean as $\pi\mu_{\mathrm{obs}}+(1-\pi)\E[B\mid R=0]$. The unobserved conditional mean lies in $[-1,1]$, giving the bounds. Any point in that interval is realized by assigning every missing response the corresponding constant score, without altering the observed response law or observation probability. This proves sharpness, including the fully missing case. The robustness criterion follows immediately from the identified set.
\end{proof}

For two treatment cells with no cross-cell restrictions, the sharp contrast interval is obtained by subtracting the upper endpoint of one cell from the lower of the other, and vice versa. Sampling uncertainty in all endpoints still requires valid inference; population bounds are not confidence intervals. Selective retention of experienced participants can otherwise imitate either learning or persistent bias.

\section{Completion Estimate}
51\%

This is a post hoc, heuristic estimate for the full catalogue target.
The attempt derives incentive and training responses in an imperfect-mental-model framework, characterizes material residual bias and changing interaction signs, and proves sharp retained-cohort attrition bounds. The full response-surface target still requires randomized stakes and feedback with measured expertise, validated task meaning, normalized payment incentives, uncertainty and attrition handling, and accurate prediction of matched consequential market decisions.

\section{Remaining gap and reference}
The effort model assumes risk neutrality in normalized payments, quadratic subjective loss, a fixed anchor and a specified feedback-encoding law. It does not identify those primitives or show that they transport to professionals. A full solution needs randomized stakes and feedback, stable task meaning, a baseline expertise definition, retained-cohort outcomes, and out-of-sample prediction for matched consequential decisions. Existing high-stakes experiments address part of that agenda; this attempt does not relabel their findings as unproved.

Benjamin Enke, Uri Gneezy, Brian Hall, David Martin, Vadim Nelidov, Theo Offerman and Jeroen van de Ven (2023), \emph{Cognitive Biases: Mistakes or Missing Stakes?}, Review of Economics and Statistics 105(4), 818--832, especially pp. 818 and 830, \url{https://doi.org/10.1162/rest_a_01093}; author manuscript \url{https://benjamin-enke.com/pdf/Biases_stakes.pdf}. The source motivates the incentive question; the mental-model calculations and attrition analysis here are the present restricted exercise.

\end{document}
